\chapter{Implémentation Logicielle et Système Temps Réel}

\section{Stack Technologique et Environnement de Développement}

La solution s'appuie sur une suite logicielle moderne, performante et standardisée dans l'industrie :
\begin{itemize}
    \item \textbf{Backend :} .NET 9, C\# 13, ASP.NET Core Web API, Entity Framework Core 9.
    \item \textbf{Médiation \& Temps Réel :} MediatR 12 (CQRS intra-processus), ASP.NET Core SignalR (WebSockets).
    \item \textbf{Persistance :} Microsoft SQL Server / LocalDB avec migrations automatiques EF Core.
    \item \textbf{Frontend :} Angular 19, TypeScript 5.6, Signals réactifs, RxJS, TailwindCSS.
    \item \textbf{Visualisation Graphique :} \texttt{vis-network} (cartographie interactive de graphe) et \texttt{Chart.js} (séries temporelles de télémétrie).
\end{itemize}

\section{Le Cœur Métier : Le DigitalTwinEngine}

Le \texttt{DigitalTwinEngine} est le service central d'orchestration backend. Il traite chaque mesure entrante de manière atomique et coordonne l'évaluation des anomalies, la régulation anti-rebond, l'analyse d'impact et la diffusion temps réel.

La Figure~\ref{fig:engine_pipeline} détaille les 7 étapes séquentielles du pipeline d'ingestion.

\begin{figure}[htbp]
\centering
\begin{tikzpicture}[
    node distance=0.8cm,
    every node/.style={font=\small\sffamily},
    stepblock/.style={
        rectangle,
        draw=ocpnavy,
        fill=ocpnavy!8,
        thick,
        rounded corners=3pt,
        text width=12cm,
        minimum height=0.9cm,
        align=left,
        inner sep=6pt
    },
    stepnum/.style={
        rectangle,
        fill=ocpnavy,
        text=white,
        font=\small\bfseries,
        rounded corners=2pt,
        inner sep=3pt
    }
]
    \node[stepblock] (s1) {\textbf{Étape 1 : Réception \& Validation Syntaxique}\\
        \scriptsize Désérialisation du DTO \texttt{IngestReadingCommand} et vérification de la présence des identifiants.};
    
    \node[stepblock, below=of s1] (s2) {\textbf{Étape 2 : Chargement du Domaine \& Validation des Invariants}\\
        \scriptsize Récupération de l'agrégat \texttt{Machine} et du \texttt{Sensor}, validation stricte via \texttt{DomainGuard}.};

    \node[stepblock, below=of s2] (s3) {\textbf{Étape 3 : Évaluation Multi-Seuils d'Anomalie}\\
        \scriptsize Analyse de la valeur par \texttt{AnomalyDetectionService} $\rightarrow$ Qualification : \textit{Normal}, \textit{Warning} ou \textit{Critical}.};

    \node[stepblock, below=of s3] (s4) {\textbf{Étape 4 : Filtrage Temporel Anti-Flood (Cooldown)}\\
        \scriptsize Contrôle par \texttt{AlertCooldownService} $\rightarrow$ Blocage des doublons si délai $<$ seuil (ex. 60s).};

    \node[stepblock, below=of s4] (s5) {\textbf{Étape 5 : Analyse d'Impact Topologique (BFS)}\\
        \scriptsize Si statut \textit{Critical}, exécution du parcours BFS sur le graphe de dépendance et calcul des machines avals impactées.};

    \node[stepblock, below=of s5] (s6) {\textbf{Étape 6 : Persistance Transactionnelle}\\
        \scriptsize Enregistrement de la mesure, mise à jour du statut machine et création de l'alerte dans SQL Server via EF Core.};

    \node[stepblock, below=of s6] (s7) {\textbf{Étape 7 : Diffusion Temps Réel WebSockets (SignalR)}\\
        \scriptsize Émission instantanée des événements \texttt{ReadingReceived}, \texttt{MachineStatusUpdated} et \texttt{AlertRaised}.};

    % Flèches de liaison
    \draw[flowarrow] (s1) -- (s2);
    \draw[flowarrow] (s2) -- (s3);
    \draw[flowarrow] (s3) -- (s4);
    \draw[flowarrow] (s4) -- (s5);
    \draw[flowarrow] (s5) -- (s6);
    \draw[flowarrow] (s6) -- (s7);
\end{tikzpicture}
\caption{Pipeline d'orchestration en 7 étapes du \texttt{DigitalTwinEngine}}
\label{fig:engine_pipeline}
\end{figure}

\section{Services Spécialisés du Backend}

\subsection{Détection d'Anomalies : AnomalyDetectionService}

Le service \texttt{AnomalyDetectionService} évalue la dérive d'une mesure par rapport aux seuils configurés sur le capteur selon la règle décisionnelle :
\begin{equation}
\text{Statut}(v) = 
\begin{cases} 
\text{Critical} & \text{si } v \ge \text{Seuil}_{\text{Critical}} \\
\text{Warning} & \text{si } \text{Seuil}_{\text{Warning}} \le v < \text{Seuil}_{\text{Critical}} \\
\text{Normal} & \text{si } v < \text{Seuil}_{\text{Warning}}
\end{cases}
\end{equation}

\subsection{Régulation Anti-Rebond : AlertCooldownService}

Pour éliminer l'inondation d'alarmes, le \texttt{AlertCooldownService} implémente un dictionnaire concurrent thread-safe (\texttt{ConcurrentDictionary<Guid, DateTime>}) qui stocke l'horodatage de la dernière alerte émise par capteur :

\begin{lstlisting}[language=CSharpCustom, caption={Implémentation du service anti-rebond AlertCooldownService}]
public class AlertCooldownService : IAlertCooldownService
{
    private readonly ConcurrentDictionary<Guid, DateTime> _lastAlertTimestamps = new();
    private readonly TimeSpan _cooldownDuration = TimeSpan.FromSeconds(60);

    public bool ShouldEmitAlert(Guid sensorId, DateTime currentTimestamp)
    {
        if (_lastAlertTimestamps.TryGetValue(sensorId, out var lastTime))
        {
            if (currentTimestamp - lastTime < _cooldownDuration)
                return false; // Dans la période de refroidissement : filtrage de l'alerte
        }

        _lastAlertTimestamps[sensorId] = currentTimestamp;
        return true; // Délai écoulé : émission autorisée
    }
}
\end{lstlisting}

\section{Architecture de Communication Temps Réel avec SignalR}

La communication entre le backend et l'interface utilisateur s'effectue via un hub WebSockets ASP.NET Core SignalR (\texttt{DigitalTwinHub}). Les clients souscrivent à des canaux logiques ciblés afin d'optimiser la bande passante réseau :

\begin{itemize}
    \item \textbf{Groupe Global (\texttt{"AllClients"}) :} Reçoit les notifications d'alertes majeures et les changements globaux de statut de l'usine.
    \item \textbf{Groupe Par Machine (\texttt{"Machine\_\{id\}"}) :} Reçoit le flux haute fréquence des mesures de télémétrie d'un équipement sélectionné par l'opérateur.
\end{itemize}

\section{Le Simulateur Industriel Autonome}

Le simulateur d'usine est développé sous la forme d'un \textbf{Worker Service} indépendant. Il reproduit fidèlement la dynamique cinétique des capteurs industriels (vibration en mm/s, température de palier en $^\circ$C, courant moteur en Ampères) et permet d'exécuter un scénario de défaillance déterministe sur le concasseur \texttt{CR-001} articulé en 6 phases :
\begin{enumerate}
    \item \textbf{Phase 1 --- Régime nominal :} Signaux stables dans la plage nominale (ex. vibration $< 3.5$\,mm/s) ;
    \item \textbf{Phase 2 --- Dégradation progressive :} Montée linéaire des vibrations franchissant le seuil Warning ($4.5$\,mm/s) ;
    \item \textbf{Phase 3 --- Défaillance critique :} Dépassement du seuil Critical ($7.0$\,mm/s), déclenchement de l'alerte critique et propagation BFS vers le convoyeur \texttt{CV-001} ;
    \item \textbf{Phase 4 --- Pic d'incident :} Maintien en zone critique simulant l'échauffement maximal ;
    \item \textbf{Phase 5 --- Refroidissement / Réparation :} Baisse progressive des valeurs de télémétrie ;
    \item \textbf{Phase 6 --- Auto-guérison et retour nominal :} Rétablissement complet de l'ensemble de la chaîne de production.
\end{enumerate}

\section{Développement de l'Interface de Supervision Angular 19}

L'application cliente est développée avec Angular 19 en exploitant le nouveau paradigme des \textbf{Signals} pour une réactivité optimale sans surcoût de cycle de détection :

\begin{itemize}
    \item \textbf{Module Dashboard :} Affiche les compteurs opérationnels globaux (machines nominales, dégradées, critiques), le journal des alertes actives et les indicateurs synthétiques de santé de l'usine.
    \item \textbf{Module Topologie de l'Usine (vis-network) :} Cartographie topologique interactive où chaque machine est représentée par un nœud dont la couleur reflète dynamiquement son statut opérationnel (Vert : Nominal, Orange : Warning, Rouge : Critical, Violet : Degraded par impact amont).
    \item \textbf{Module Détails Machine (Chart.js) :} Visualisation en temps réel des courbes chronologiques des capteurs avec tracé distinct des seuils de tolérance.
    \item \textbf{Module Gestion des Alertes :} Console de filtrage multi-critères avec interface d'acquittement instantané pour les opérateurs de salle de contrôle.
\end{itemize}
